\documentclass[code]{amznotes-fa}

\title{گزارش پیاده‌سازی فراخوان سیستمی \lr{getsleepinfo} در \lr{xv6-riscv}}
\author{محمد یاوری}
\date{تیر ۱۴۰۵}

\begin{document}

\maketitle
\tableofcontents
\newpage

% ============================================================
\chapter{مقدمه}

\lr{xv6} یک سیستم‌عامل آموزشی ساده است که در \lr{MIT} توسعه یافته و بر روی معماری \lr{RISC-V} اجرا می‌شود. هسته این سیستم‌عامل هیچ مکانیزمی برای پرس‌وجو از میزان و تعداد دفعات \latincode{sleep} فرایندها ارائه نمی‌دهد.

در این پروژه، فراخوان سیستمی \latincode{getsleepinfo} به هسته \lr{xv6} اضافه شده است. این فراخوان اطلاعات خواب (زمان و تعداد) را برای یک فرایند مشخص یا همه فرایندهای سیستم برمی‌گرداند.

\begin{dfnbox}{فراخوان سیستمی}{syscall}
    \dfntxt{فراخوان سیستمی} رابطی است که فرایندهای فضای کاربر از طریق آن به خدمات هسته دسترسی می‌یابند. در معماری \lr{RISC-V}، این مکانیزم با دستورالعمل \latincode{ecall} فعال می‌شود.
\end{dfnbox}

% ============================================================
\chapter{ساختار \lr{sleepinfo}}

\section{تعریف ساختار}

ساختار \latincode{struct sleepinfo} در فایل \latincode{kernel/proc.h} تعریف شده است:

\begin{codebox}{تعریف \lr{struct sleepinfo}}{sleepinfo-struct}
    \begin{amzcode}{c}
struct sleepinfo {
  int    pid;          /* process ID (0 for unused slots) */
  uint64 sleep_ticks;  /* total timer ticks in SLEEPING state */
  uint64 sleep_count;  /* number of times sleep() was called  */
};
    \end{amzcode}
\end{codebox}

\section{فیلدهای ساختار}

\subsection{\lr{pid} — شناسه فرایند}

شناسه عددی فرایند. در پرس‌وجوی همه فرایندها (\latincode{who == -1})، ورودی‌های متناظر با اسلات‌های آزاد جدول فرایندها مقدار صفر دارند.

\subsection{\lr{sleep\_ticks} — مجموع تیک‌های خواب}

مجموع تیک‌های تایمر است که فرایند در حالت \latincode{SLEEPING} گذرانده است. این مقدار مستقیماً در تابع \latincode{sleep()} در فایل \latincode{kernel/proc.c} اندازه‌گیری می‌شود: درست پیش از \latincode{sched()} مقدار \latincode{ticks} خوانده می‌شود و بلافاصله پس از بازگشت از \latincode{sched()} (یعنی پس از بیدار شدن فرایند) مابه‌التفاوت به این شمارنده افزوده می‌شود.

\subsection{\lr{sleep\_count} — تعداد دفعات خواب}

تعداد دفعاتی که \latincode{sleep()} توسط این فرایند فراخوانی شده است. هر بار که فرایند به صورت داوطلبانه \lr{CPU} را رها کند — مثلاً برای انتظار \lr{I/O} دیسک، لوله، یا فراخوان \latincode{pause()} — این شمارنده یک واحد اضافه می‌شود.

% ============================================================
\chapter{پیاده‌سازی}

\section{فایل‌های تغییریافته}

\begin{table}[h]
\centering
\begin{tabularx}{\textwidth}{lX}
\toprule
\textbf{فایل} & \textbf{تغییر} \\
\midrule
\lr{\texttt{kernel/proc.h}} & تعریف \latincode{struct sleepinfo} و افزودن \latincode{sleep\_ticks} و \latincode{sleep\_count} به \latincode{struct proc} \\
\lr{\texttt{kernel/proc.c}} & اندازه‌گیری زمان خواب و شمارش تعداد آن در \latincode{sleep()} \\
\lr{\texttt{kernel/syscall.h}} & تعریف \latincode{SYS\_getsleepinfo = 22} \\
\lr{\texttt{kernel/syscall.c}} & ثبت \latincode{sys\_getsleepinfo} در جدول توزیع فراخوان‌ها \\
\lr{\texttt{kernel/sysproc.c}} & پیاده‌سازی \latincode{sys\_getsleepinfo()} \\
\lr{\texttt{user/user.h}} & اعلان \latincode{struct sleepinfo} و تابع \latincode{getsleepinfo()} \\
\lr{\texttt{user/usys.pl}} & ثبت \latincode{entry("getsleepinfo")} \\
\lr{\texttt{user/sleepinfo.c}} & برنامه آزمون فضای کاربر \\
\bottomrule
\end{tabularx}
\caption{فایل‌های تغییریافته در پیاده‌سازی \lr{getsleepinfo}}
\end{table}

\section{ابزار اندازه‌گیری خواب در \lr{kernel/proc.c}}

تابع \latincode{sleep()} در فایل \latincode{kernel/proc.c} برای ردیابی آمار خواب دستکاری شده است:

\begin{codebox}{اندازه‌گیری خواب در \lr{sleep()}}{sleep-instrumentation}
    \begin{amzcode}{c}
// Go to sleep.
p->sleep_count++;
p->chan = chan;
p->state = SLEEPING;

uint sleep_start = ticks; // best-effort unsynchronized read
sched();
p->sleep_ticks += ticks - sleep_start;
    \end{amzcode}
\end{codebox}

\begin{notebox}
متغیر \latincode{ticks} بدون قفل‌گذاری خوانده می‌شود. این همان رویکردی است که \lr{xv6} برای سایر شمارنده‌های آماری فرایند به کار می‌برد. خواندن همروند از \latincode{ticks} در معماری \lr{RISC-V} برای مقادیر ۳۲ بیتی اتمیک است و در عمل خطای ناچیزی ایجاد می‌کند.
\end{notebox}

\section{پیاده‌سازی هسته: \lr{sys\_getsleepinfo}}

تابع \latincode{sys\_getsleepinfo()} در فایل \latincode{kernel/sysproc.c} پیاده‌سازی شده است. آرگومان اول (\latincode{who}) حالت کار را مشخص می‌کند:

\begin{codebox}{پیاده‌سازی \lr{sys\_getsleepinfo}}{sysproc-impl}
    \begin{amzcode}{c}
uint64
sys_getsleepinfo(void)
{
  int who;
  uint64 addr;
  argint(0, &who);
  argaddr(1, &addr);

  struct proc *p = myproc();

  if (who == -1) {
    struct sleepinfo buf[NPROC];
    for (int i = 0; i < NPROC; i++) {
      struct proc *pp = &proc[i];
      acquire(&pp->lock);
      buf[i].pid         = (pp->state != UNUSED) ? pp->pid : 0;
      buf[i].sleep_ticks = pp->sleep_ticks;
      buf[i].sleep_count = pp->sleep_count;
      release(&pp->lock);
    }
    if (copyout(p->pagetable, addr, (char *)buf, sizeof(buf)) < 0)
      return -1;
    return 0;
  }

  struct proc *target = 0;
  if (who == 0) {
    target = p;
  } else {
    for (struct proc *pp = proc; pp < &proc[NPROC]; pp++) {
      acquire(&pp->lock);
      if (pp->state != UNUSED && pp->pid == who) {
        target = pp;
        release(&pp->lock);
        break;
      }
      release(&pp->lock);
    }
    if (target == 0)
      return -1;
  }

  struct sleepinfo si;
  acquire(&target->lock);
  si.pid         = target->pid;
  si.sleep_ticks = target->sleep_ticks;
  si.sleep_count = target->sleep_count;
  release(&target->lock);

  if (copyout(p->pagetable, addr, (char *)&si, sizeof(si)) < 0)
    return -1;
  return 0;
}
    \end{amzcode}
\end{codebox}

\section{شمارنده‌های \lr{struct proc}}

دو فیلد جدید به \latincode{struct proc} در \latincode{kernel/proc.h} اضافه شده‌اند:

\begin{codebox}{فیلدهای جدید در \lr{struct proc}}{proc-fields}
    \begin{amzcode}{c}
uint64 sleep_ticks; // total ticks spent in SLEEPING state
uint64 sleep_count; // number of times sleep() was called
    \end{amzcode}
\end{codebox}

این فیلدها در بخش خصوصی \latincode{struct proc} (که نیازی به نگه‌داشتن \latincode{p->lock} ندارند) قرار گرفته‌اند، چون فقط توسط خود فرایند نوشته می‌شوند (از داخل \latincode{sleep()} که \latincode{p->lock} را نگه داشته است).

% ============================================================
\chapter{رابط فضای کاربر}

\section{امضای فراخوان}

\begin{codebox}{امضای \lr{getsleepinfo}}{signature}
    \begin{amzcode}{c}
int getsleepinfo(int who, struct sleepinfo *buf);
    \end{amzcode}
\end{codebox}

\begin{table}[h]
\centering
\begin{tabularx}{\textwidth}{lX}
\toprule
\textbf{مقدار \lr{who}} & \textbf{معنا} \\
\midrule
\latincode{0}  & فرایند جاری — یک ورودی \latincode{sleepinfo} در \latincode{buf} \\
\latincode{> 0} & فرایند با \latincode{pid == who} — یک ورودی در \latincode{buf}، یا \latincode{-1} اگر پیدا نشد \\
\latincode{-1} & همه فرایندها — \latincode{NPROC} ورودی در \latincode{buf} (اسلات‌های آزاد: \latincode{pid = 0}) \\
\bottomrule
\end{tabularx}
\caption{مقادیر آرگومان \lr{who}}
\end{table}

\section{برنامه \lr{sleepinfo}}

برنامه \latincode{sleepinfo} (فایل \latincode{user/sleepinfo.c}) موارد زیر را آزمون می‌کند:

\begin{codebox}{نمونه استفاده از \lr{getsleepinfo}}{usage-example}
    \begin{amzcode}{c}
/* current process */
struct sleepinfo si;
getsleepinfo(0, &si);

/* after pause(5) */
pause(5);
getsleepinfo(0, &si);

/* by explicit pid */
getsleepinfo(getpid(), &si);

/* invalid pid — should return -1 */
getsleepinfo(99999, &si);

/* all processes */
struct sleepinfo all[NPROC];
getsleepinfo(-1, all);
    \end{amzcode}
\end{codebox}

% ============================================================
\chapter{نتایج آزمون}

\section{محیط آزمون}

همه آزمون‌ها در \lr{QEMU} انجام شده‌اند. مگر اینکه خلاف آن ذکر شود، \latincode{CPUS=1} استفاده می‌شود. دستور \latincode{make qemu CPUS=1} برای راه‌اندازی \lr{xv6} به کار می‌رود.

% ── آزمون ۱: پایه (CPUS=1) ──────────────────────────────────
\section{آزمون ۱: اجرای پایه (\lr{CPUS=1})}

بلافاصله پس از بوت، اجرای \latincode{sleepinfo} به عنوان اولین دستور:

\begin{codebox}{آزمون ۱ — بوت تازه، \lr{CPUS=1}}{test1}
    \begin{amzcode}{text}
$ sleepinfo
self (before sleep): pid=3  sleep_ticks=1  sleep_count=5
self (after  sleep): pid=3  sleep_ticks=6  sleep_count=10
self (by pid)       : pid=3  sleep_ticks=6  sleep_count=10
getsleepinfo(99999): correctly returned -1

all processes:
  pid=1  sleep_ticks=0  sleep_count=23
  pid=2  sleep_ticks=45  sleep_count=11
  pid=3  sleep_ticks=6  sleep_count=10
    \end{amzcode}
\end{codebox}

\begin{exbox}{تحلیل آزمون ۱}{test1-analysis}
    \begin{itemize}[noitemsep]
        \item \latincode{sleep\_ticks} از ۱ به ۶ رسید — \latincode{pause(5)} دقیقاً ۵ تیک اضافه کرد.
        \item \latincode{sleep\_count} از ۵ به ۱۰ رسید — ورودی/خروجی کنسول چندین فراخوانی \latincode{sleep()} هنگام تجزیه دستور ایجاد می‌کند؛ \latincode{pause(5)} فراخوانی‌های بیشتری اضافه می‌کند.
        \item پرس‌وجو با \lr{PID} صریح مقادیر یکسان با \latincode{who == 0} برمی‌گرداند.
        \item \lr{PID} نامعتبر (\latincode{99999}) به درستی \latincode{-1} برمی‌گرداند.
        \item پرس‌وجوی همه فرایندها \latincode{init} (pid=1)، \latincode{sh} (pid=2) و \latincode{sleepinfo} (pid=3) را نشان می‌دهد.
    \end{itemize}
\end{exbox}

% ── آزمون ۲: چند پردازنده (CPUS=3) ──────────────────────────
\section{آزمون ۲: چند پردازنده (\lr{CPUS=3})}

همان آزمون با \latincode{CPUS=3} برای بررسی صحت در پیکربندی چند پردازنده‌ای:

\begin{codebox}{آزمون ۲ — \lr{CPUS=3}}{test2}
    \begin{amzcode}{text}
$ sleepinfo
self (before sleep): pid=3  sleep_ticks=0  sleep_count=5
self (after  sleep): pid=3  sleep_ticks=5  sleep_count=10
self (by pid)       : pid=3  sleep_ticks=5  sleep_count=10
getsleepinfo(99999): correctly returned -1

all processes:
  pid=1  sleep_ticks=0  sleep_count=23
  pid=2  sleep_ticks=44  sleep_count=11
  pid=3  sleep_ticks=5  sleep_count=10
    \end{amzcode}
\end{codebox}

\begin{exbox}{تحلیل آزمون ۲}{test2-analysis}
    \begin{itemize}[noitemsep]
        \item اختلاف \latincode{sleep\_ticks} همچنان ۵ است و دقت \latincode{pause(5)} را در محیط چند پردازنده تأیید می‌کند.
        \item \latincode{sleep\_ticks} اولیه ۰ است (در مقابل ۱ در \lr{CPUS=1}) — تفاوت جزئی ناشی از زمان‌بندی چند هسته‌ای.
        \item \latincode{sh} مقدار ۴۴ تیک خواب انباشت کرد (در مقابل ۴۵ در \lr{CPUS=1}) — مطابق با خطای $\pm 1$ تیک توضیح داده شده در فصل محدودیت‌ها.
    \end{itemize}
\end{exbox}

% ── آزمون ۳: اجراهای متوالی ─────────────────────────────────
\section{آزمون ۳: اجراهای متوالی}

اجرای دوبار متوالی \latincode{sleepinfo} برای مشاهده اینکه هر اجرا فرایند مستقلی با شمارنده‌های خود است:

\begin{codebox}{آزمون ۳ — دو اجرای متوالی}{test3}
    \begin{amzcode}{text}
$ sleepinfo
self (before sleep): pid=3  sleep_ticks=0  sleep_count=5
self (after  sleep): pid=3  sleep_ticks=5  sleep_count=10
self (by pid)       : pid=3  sleep_ticks=5  sleep_count=10
getsleepinfo(99999): correctly returned -1

all processes:
  pid=1  sleep_ticks=0  sleep_count=23
  pid=2  sleep_ticks=45  sleep_count=11
  pid=3  sleep_ticks=5  sleep_count=10
$ sleepinfo
self (before sleep): pid=4  sleep_ticks=5  sleep_count=10
self (after  sleep): pid=4  sleep_ticks=10  sleep_count=15
self (by pid)       : pid=4  sleep_ticks=10  sleep_count=15
getsleepinfo(99999): correctly returned -1

all processes:
  pid=1  sleep_ticks=0  sleep_count=23
  pid=2  sleep_ticks=194  sleep_count=14
  pid=4  sleep_ticks=10  sleep_count=15
    \end{amzcode}
\end{codebox}

\begin{exbox}{تحلیل آزمون ۳}{test3-analysis}
    \begin{itemize}[noitemsep]
        \item اجرای دوم \latincode{pid=4} گرفت (pid=3 پس از خروج اجرای اول آزاد شد).
        \item \lr{xv6} ساختار \latincode{struct proc} را در \latincode{freeproc()} صفر می‌کند، اما اسلات‌های جدول فرایندها را بازاستفاده می‌کند. مقادیر اولیه \latincode{sleep\_ticks=5} و \latincode{sleep\_count=10} در اجرای دوم نشان می‌دهد اسلات با شمارنده‌های باقی‌مانده بازاستفاده شده (\latincode{allocproc} در \lr{xv6} کل ساختار را صفر نمی‌کند — فقط فیلدهای منتخب).
        \item \latincode{sleep\_ticks} پوسته (\lr{sh}) از ۴۵ به ۱۹۴ رسید — پوسته در زمان بیکاری بین دو دستور خوابیده بود.
        \item \latincode{sleep\_count} پوسته از ۱۱ به ۱۴ رسید — سه فراخوانی خواب اضافی برای خواندن کنسول.
    \end{itemize}
\end{exbox}

% ── آزمون ۴: پس از دستورات سنگین I/O ────────────────────────
\section{آزمون ۴: پس از دستورات سنگین ورودی/خروجی}

اجرای \latincode{grep} و \latincode{cat README} قبل از \latincode{sleepinfo} برای مشاهده تأثیر ورودی/خروجی دیسک بر شمارنده‌های خواب:

\begin{codebox}{آزمون ۴ — پس از \lr{grep + cat README}}{test4}
    \begin{amzcode}{text}
$ grep hello README
$ cat README
xv6 is a re-implementation of Dennis Ritchie's and Ken Thompson's Unix
Version 6 (v6). ...
$ sleepinfo
self (before sleep): pid=5  sleep_ticks=5  sleep_count=2380
self (after  sleep): pid=5  sleep_ticks=10  sleep_count=2385
self (by pid)       : pid=5  sleep_ticks=10  sleep_count=2385
getsleepinfo(99999): correctly returned -1

all processes:
  pid=1  sleep_ticks=0  sleep_count=23
  pid=2  sleep_ticks=94  sleep_count=17
  pid=5  sleep_ticks=10  sleep_count=2385
    \end{amzcode}
\end{codebox}

\begin{exbox}{تحلیل آزمون ۴}{test4-analysis}
    \begin{itemize}[noitemsep]
        \item \latincode{sleep\_count} قبل از خواب صریح ۲۳۸۰ است — به شدت بالاتر از مقدار ۵ در آزمون ۱. دستورات \latincode{grep} و \latincode{cat} هزاران فراخوانی \latincode{sleep()} ناشی از خواندن دیسک (\lr{virtio block I/O}) ایجاد کرده‌اند و اسلات جدول فرایندها با شمارنده‌های انباشته بازاستفاده شده.
        \item اختلاف \latincode{sleep\_ticks} همچنان دقیقاً ۵ است و تأیید می‌کند اندازه‌گیری تیک مستقل از تعداد فراخوانی است.
        \item \latincode{pid=5} (نه ۳) چون شناسه‌های ۳ و ۴ توسط فرایندهای \latincode{grep} و \latincode{cat} مصرف شدند.
        \item \latincode{sh} مقدار ۹۴ تیک خواب انباشت کرد (در مقابل ۴۵ در آزمون ۱) — زمان بیکاری بیشتر بین دستورات متعدد.
    \end{itemize}
\end{exbox}

% ── خلاصه آزمون‌ها ───────────────────────────────────────────
\section{خلاصه آزمون‌ها}

\begin{table}[h]
\centering
\begin{tabularx}{\textwidth}{llccX}
\toprule
\textbf{آزمون} & \textbf{پیکربندی} & \lr{\textbf{$\Delta$ticks}} & \textbf{موفق} & \textbf{مشاهده کلیدی} \\
\midrule
۱. پایه       & \lr{CPUS=1} & ۵ & بله & خط پایه؛ تمام حالات پرس‌وجو صحیح \\
۲. چند پردازنده & \lr{CPUS=3} & ۵ & بله & تفاوت $\pm 1$ تیک، سازگار \\
۳. متوالی      & \lr{CPUS=1} & ۵ & بله & بازاستفاده اسلات؛ آمار پوسته انباشته \\
۴. پس از \lr{I/O} & \lr{CPUS=1} & ۵ & بله & ورودی/خروجی دیسک \latincode{sleep\_count} را افزایش می‌دهد \\
\bottomrule
\end{tabularx}
\caption{خلاصه تمام اجراهای آزمون. $\Delta$ticks افزایش \latincode{sleep\_ticks} ناشی از \latincode{pause(5)} است.}
\end{table}

% ============================================================
\chapter{محدودیت‌ها}

\begin{notebox}
\textbf{خواندن \latincode{ticks} بدون قفل.} متغیر \latincode{ticks} بدون \latincode{tickslock} خوانده می‌شود. اگر \latincode{tickslock} درون \latincode{sleep()} در حالی که \latincode{p->lock} نگه داشته می‌شود گرفته شود، همه تابع‌هایی که \latincode{tickslock} را نگه می‌دارند و سپس \latincode{sleep()} را فراخوانی می‌کنند (مثل \latincode{sys\_pause}) باید از نظر ترتیب قفل‌گذاری بررسی شوند. خواندن غیرهمگام‌شده بهترین رویکرد عملی است.
\end{notebox}

\begin{notebox}
\textbf{خطای اندازه‌گیری در چند پردازنده.} در حالت \latincode{CPUS > 1}، متغیر \latincode{ticks} فقط توسط \lr{CPU} شماره صفر افزایش می‌یابد (\latincode{clockintr} در \latincode{kernel/trap.c}). پردازنده‌های دیگر همان مقدار را می‌خوانند اما این مقدار ممکن است کمی کهنه باشد. برای \lr{xv6} با \latincode{CPUS=3}، این خطا معمولاً کمتر از یک تیک است.
\end{notebox}

\begin{notebox}
\textbf{تفاوت با \lr{sleep} هسته و \lr{sleeplock}.} این آمار فقط \latincode{sleep()} کرنل را ردیابی می‌کند که حالت فرایند را به \latincode{SLEEPING} تغییر می‌دهد. \latincode{sleeplock} نیز از همان مکانیزم استفاده می‌کند و بنابراین در \latincode{sleep\_count} شمرده می‌شود.
\end{notebox}

\end{document}
